% Generated draft: requires mathematical and source review.
\documentclass[11pt]{article}
\usepackage[margin=25mm]{geometry}
\usepackage{fontspec}
\setmainfont{Latin Modern Roman}
\usepackage{amsmath,amssymb,mathtools,unicode-math}
\setmathfont{Latin Modern Math}
\usepackage{enumitem,needspace,xurl,hyperref}
\setlength{\parindent}{0pt}
\setlength{\parskip}{5pt}
\setlength{\emergencystretch}{3em}
\newcommand{\N}{\mathbb{N}}
\newcommand{\Z}{\mathbb{Z}}
\newcommand{\Q}{\mathbb{Q}}
\newcommand{\R}{\mathbb{R}}
\newcommand{\C}{\mathbb{C}}
\newcommand{\F}{\mathbb{F}}
\newcommand{\A}{\mathbb{A}}
\newcommand{\PP}{\mathbb{P}}
\newcommand{\E}{\mathbb{E}}
\newcommand{\eps}{\varepsilon}
\newcommand{\dd}{\,\mathrm{d}}
\begin{document}
\section*{UnsolvedMath ID: 20000240}
\textbf{Source problem code:} catalogue entry
\subsection{Definitions and mathematical statement}
The supplied reconstruction of Christian Haesemeyer's question is: ``What are motivic local systems? Can one define motivic Atiyah classes?'' The original extraction is damaged. The question asks for definitions and a construction, not a theorem about an already specified category. The report distinguishes $\A^1$-homotopy, Nori, stable-motivic and geometric-origin meanings of local system; they are not silently identified. It claims a precise rational characteristic-zero answer in a perverse Nori framework, without giving that framework's definition or reference, and says no standard intrinsic construction called a motivic Atiyah class was located. This is a limited supplied literature assessment, not a nonexistence theorem.

For comparison, if $\mathcal V$ is a vector bundle on a smooth characteristic-zero variety $X$, its ordinary coherent Atiyah class is the extension class
\[
 \operatorname{At}(\mathcal V)\in
 \operatorname{Ext}^1_X(\mathcal V,\mathcal V\otimes\Omega_X^1)
\]
of the first-jet sequence $0\to\mathcal V\otimes\Omega_X^1\to J^1\mathcal V\to\mathcal V\to0$. Vanishing means that this sequence admits a splitting, equivalently an algebraic connection; integrability is an additional property. This coherent invariant, called ordinary de Rham Atiyah class in the report, is not by definition an intrinsic motivic Atiyah class.

The report claims that for every finite etale $p:Y\to X$ over a smooth characteristic-zero base, $A=p_*\mathcal O_Y$ has the canonical integrable connection induced by $d_Y$ and the etale identification $\Omega_Y^1\cong p^*\Omega_X^1$. Its trace-zero summand $A_0=\ker(\operatorname{Tr}:A\to\mathcal O_X)$ is horizontal and has vanishing Atiyah class, as does $A$. For a constant positive degree on each connected component, division by that degree gives the stated trace splitting in characteristic zero.

More specifically let $E/\C$ be any elliptic curve and let $p=[m]:E\to E$ for any integer $m\geq2$. The trace-zero Artin local system and its underlying algebraic bundle are
\[
 \mathcal L_0=\ker([m]_*\Q_E\xrightarrow{\operatorname{Tr}}\Q_E),\qquad
 A_0=\ker([m]_*\mathcal O_E\xrightarrow{\operatorname{Tr}}\mathcal O_E).
\]
The report claims that $\mathcal L_0$ has nontrivial finite monodromy, while
\[
 \operatorname{rank} A_0=m^2-1,\qquad H^0(E,A_0)=0,\qquad
 A_0\not\cong\mathcal O_E^{\oplus(m^2-1)},\qquad\operatorname{At}(A_0)=0.
\]
The full pushforward has regular deck-group monodromy and the trace-zero part has its augmentation representation. Consequently this ordinary Atiyah-class test cannot detect triviality or monodromy even on these trace-free Artin examples. The example does not exclude a richer motivic class with different target or information.
\subsection{Short English statement}
Define motivic local systems and seek an intrinsic motivic Atiyah class, keeping the possible motivic categories distinct. The report supplies finite-etale examples with flat connections: trace-free local systems from multiplication on an elliptic curve retain nontrivial monodromy although their underlying bundles have zero ordinary Atiyah class. This limits that ordinary invariant, rather than resolving every motivic definition.
\subsection{Sources}
\begin{itemize}
\item Local registry: \url{data/unsolvedmath/problems.json}, permanent ID $20000240$, source code \texttt{catalogue entry}; damaged original text, supplied reconstruction and research summary. Statement status reconstructed and unverified.
\item AIM workshop ``Projective modules and $A^1$-homotopy theory'', Question 3, attributed to Christian Haesemeyer: \url{https://aimath.org/pastworkshops/projectiveA1problems.pdf}. Supplied verification identifies PDF page 2; summary dates the question to 2014. Canonical location \texttt{aim-algebraic-geometry-notes.json}, \texttt{notes[239]}. No PDF was opened here.
\item No original remarks or literature citation supplied. The report's perverse Nori framework and literature-search assessment have no full bibliographic references or search date.
\item Supplied machine-generated research attempt 1: partial result; problem partially solved at the run and in the registry; rigor metadata ``proved'', novelty confidence low. This non-peer-reviewed aid requires independent verification before citation.
\end{itemize}
\end{document}
